% !TEX encoding = UTF-8
% !TEX TS-program = pdflatex
% !TEX root = ../tesi.tex

%**************************************************************
\chapter{Descrizione dello stage}
\label{cap:descrizione-stage}
%**************************************************************
\section{L'azienda}
\label{azienda}
G-Squared Srl è una società di ingegneria informatica specializzata nello sviluppo di software applicativo medicale. Nata nel 2008 con sede legale a Vicenza e come spin-off della società Studio Synthesis Sistemi Avanzati srl, ora ha sede operativa a Ponte San Nicolò (PD).
\\
La filosofia di G-Squared si incentra molto sui dipendenti, concentrandosi su spirito di squadra, rispetto delle regole e valore alla persona e al merito.
\\
Uno dei punti di forza dell'azienda è l'esperienza acquisita in oltre vent'anni di attività svolta in ambito medicale. Questo permette a G-Squared di offrire ai propri clienti una serie di servizi, tra cui:
\begin{itemize}
	\item Interfacciamento con macchine diagnostiche digitali;
	\item Consulenze per reti di macchine diagnostiche;
	\item Elaborazione di immagini diagnostiche specializzate.
\end{itemize}

Lavorando con software medici, come visualizzatori \textit{DICOM}, G-Squared è un'azienda certificata ISO 13485:2016. Possiede inoltre molte certificazioni per i propri software, che nel caso delle immagini DICOM si attengono alle relative dichiarazioni di conformità. Più in generale, molti dei software prodotti da G-Squared sono stati classificati come dispositivo di classe I secondo i criteri della Direttiva 1993/42/CE.

\begin{figure}[ht]
    \centering
    \includegraphics[width=0.8\textwidth]{immagini/logo-azienda.png}
    \caption{\textit{Logo azienda G-Squared}}
    \textbf{Fonte}: \href{https://www.gsquared.it/it}{gsquared.it}
    \label{fig: Logo azienda G-Squared}
\end{figure}
\newpage

%**************************************************************
\section{Scelta dello stage}\label{sec:scelta-dello-stage}
Da anni ho iniziato ad interessarmi di più al rendering real-time e alle sue varie applicazioni: dai videogiochi, al training, fino a simulazioni di vario genere.
Ho partecipato all'evento Stage-IT 2019, ma durante la ricerca di uno stage a fine 2019 tra i vari colloqui ho deciso di indagare più in profondità, cercando le opportunità presentate agli eventi Stage-IT degli anni passati.
\\
Dopo varie ricerche, ho trovato G-Squared nel documento di Stage-IT 2017. Contattando e confrontandomi con l'azienda, ho deciso di unire la mia passione per C\texttt{++} e per il rendering a un ambito utile, come quello medicale, in modo da imparare di più su come funziona e tutte le tecnologie che vengono impiegate.

%**************************************************************
\section{Introduzione al progetto e interessi aziendali}
G-Squared attraverso questo progetto mira ad analizzare la possibilità di integrare un visualizzatore volumetrico in uno dei suoi software. Se fatto correttamente questo può essere vantaggioso: permettendo al medico di analizzare il volume tridimensionale insieme alle immagini radiologiche, permettendo così una visione più completa.

%**************************************************************
\section{Obiettivi}\label{sec:descrizione-obiettivi}
L'obiettivo di questo stage è lo sviluppo di un widget Qt che permetta la visualizzazione 3D volumetrica di una ricostruzione fatta tramite immagini diagnostiche medicali radiologiche. Dovevo quindi analizzare le funzionalità del framework VTK e la sua integrazione con il framework Qt, per sviluppare un widget che permettesse di caricare un volume da un esame radiologico e di visualizzarlo. Era richiesta la possibilità di applicarci dei filtri tramite dei Preset prestabiliti e la possibilità di tagliare il volume.
Importante era la compatibilità con le librerie aziendali, in modo da rendere il widget facilmente integrabile con gli strumenti e i database di gestione delle immagini già presenti e utilizzati in azienda.
\\
Durante la stesura del piano di lavoro con il tutor aziendale, che è avvenuta prima dell'inizio dello stage, sono stati individuati obiettivi suddivisi in obbligatori e desiderabili.

\subsection{Obiettivi obbligatori}\label{sec:obiettivi-obbligatori}
\begin{itemize}
\item \underline{\textit{O01}}: visualizzazione interattiva volumetrica;
\item \underline{\textit{O02}}: possibilità di scelta della funzione di trasferimento dei voxel (colore, trasparenza).
\end{itemize}

\subsection{Obiettivi desiderabili}\label{sec:obiettivi-desiderabili}
\begin{itemize}
\item \underline{\textit{D01}}: piani di taglio del volume;
\item \underline{\textit{D02}}: modifiche alla funzione taglio;
\item \underline{\textit{D03}}: ottimizzazione rendering GPU.
\end{itemize}

\subsection{Obiettivi facoltativi}
\begin{itemize}
	 \item \underline{\textit{F01}}: algoritmi di segmentazione con ITK;
	 \item \underline{\textit{F02}}: analisi unit-testing su GUI-Qt;
	 \item \underline{\textit{F03}}: porting librerie aziendali su CMake.
\end{itemize}

%**************************************************************
\section{Contesto Applicativo}
G-Squared Srl offre vari software e servizi per elaborare, archiviare e distribuire immagini radiologiche. Uno dei più importanti per il mio stage è EyeRad.

\begin{figure}[h]
    \centering
    \includegraphics[width=0.5\textwidth]{immagini/logo-software-eyerad.png}
    \caption{\textit{Logo EyeRad G-Squared}}
    \textbf{Fonte}: \href{https://www.gsquared.it/it/pacs-ris-web-gateway-dicom-ihe-2/software-immagini-medicali/12-eyerad-refertazione-digitale}{gsquared.it/12-eyerad-refertazione-digitale}
    \label{fig: Logo software EyeRad G-Squared}
\end{figure}

EyeRad è un sistema di acquisizione, analisi e gestione di immagini digitali provenienti da sistemi che aderiscono allo standard DICOM in ambiente Microsoft Windows.
EyeRad è in grado di visualizzare, elaborare, archiviare, inviare e stampare immagini digitali.
\\
\'E nato nato per soddisfare le esigenze di efficienza e precisione richieste ad una workstation per la refertazione radiologica: avvalendosi dell'uso di monitor LCD ad alta risoluzione, 2-3-5 MPixel e sfruttando appieno la possibilità di tali monitor di visualizzare le immagini digitali con 2048 toni di grigio.
\\
Nel caso del mio stage non dovevo modificarlo direttamente, ma è stato molto interessante vedere come funziona, com'è strutturato e com'è mantenuto quando sono necessarie delle modifiche o delle correzioni.

%**************************************************************
\section{Tecnologie Utilizzate}
\subsection{C\texttt{++}}\label{sec:C++}
C\texttt{++} è un linguaggio di programmazione general-purpose, sviluppato come evoluzione del linguaggio C inserendo la programmazione orientata agli oggetti. Col tempo ha avuto notevoli evoluzioni, come l'introduzione dell'astrazione rispetto al tipo. \'E il linguaggio principale usato da G-Squared per i software desktop, di conseguenza la sua scelta è stata praticamente obbligatoria. Si è deciso inoltre di utilizzare lo standard C\texttt{++}17, in modo da poter usufruire, se necessario, di tutte le novità presenti nel linguaggio come le lambda-espressioni.

\subsection{Qt5}\label{sec:Qt5}
Qt è un \textit{framework} open-source multipiattaforma per lo sviluppo di programmi con interfaccia grafica tramite l'uso di \textit{widget}.
\'E il framework principale utilizzato da G-Squared con C\texttt{++} per sviluppare interfacce grafiche, ed è stato quindi il framework che ho utilizzato per sviluppare il widget.

%**************************************************************
\section{Strumenti Utilizzati}
\subsection{CMake}\label{sec:cmake}
CMake è un software open-source multipiattaforma per l'automazione dello sviluppo, il cui nome è un'abbreviazione di "cross platform make". Questo software nasce per rimpiazzare Automake nella generazione dei Makefile, cercando di essere più semplice da usare.
\'E stato scelto per dimostrare come poter compilare e gestire un progetto C\texttt{++} con Qt, sostituendo il sistema di generazione proprietario di Qt chiamato qmake.

\subsection{git}
Git è un software di controllo di versione distribuito, utilizzato per fare il versionamento e l'upload del codice nel repository aziendale. Non sono state utilizzate regole particolari come i feature branch. 

\subsection{Visual Studio}\label{sec:visual-studio}
Microsoft Visual Studio, chiamato più comunemente Visual Studio, è un ambiente di sviluppo integrato (Integrated development environment o IDE) sviluppato da Microsoft.
Visual Studio è multi-linguaggio, ma è stato utilizzato con C\texttt{++} con il compilatore di Microsoft chiamato MSVC.

\subsection{Qt Creator}\label{sec:qt-creator}
Qt Creator è un è un IDE open-source, che semplifica lo sviluppo di applicazioni con una GUI. Fa parte dell'SDK di Qt ed è ben integrato con tutti i relativi strumenti. Permette la scelta del compilatore, che nel mio caso è stato MSVC, in modo da utilizzare lo stesso compilatore tra Visual Studio e Qt Creator.
\\
\'E il principale IDE utilizzato da G-Squared, quindi era importante verificare che tutto funzionasse anche su Qt Creator.
\\
Entrambi gli IDE: sia Visual Studio che Qt Creator permettono di aprire direttamente un progetto Cmake da un CMakeLists.txt.

%**************************************************************
\section{Aspettative personali}
Come descritto nella sezione \nameref{sec:scelta-dello-stage} (§\ref{sec:scelta-dello-stage}), ero alla ricerca di uno stage che mi permettesse di esplorare meglio gli ambiti e gli utilizzi della grafica real-time. Inoltre, mi sarebbe piaciuto utilizzare e migliorare le mie conoscenze di C\texttt{++}. Proprio per questi motivi, tra i vari colloqui effettuati, la mia scelta è ricaduta su G-Squared. I miei obiettivi da raggiungere in questo progetto di stage quindi erano:
\begin{itemize}
\item imparare come vengono archiviate e gestite le immagini mediche e gli standard utilizzati, come DICOM;
\item imparare le basi del volume rendering utilizzato in ambito medico;
\item migliorare le mie conoscenze e abilità nell'utilizzare C\texttt{++} e gli strumenti correlati.
\end{itemize}